\begin{abstract}
\textbf{Background:} Dermatological deep learning models have been shown to exhibit performance disparities across skin tones, with diagnostic accuracy often degrading substantially for darker skin presentations \cite{Daneshjou2022}. Such biases threaten the equitable deployment of clinical decision support tools and may exacerbate existing health disparities in dermatology.

\textbf{Methods:} We evaluate melanoma classification bias on the Dermatology Data Initiative (DDI) dataset \cite{Tschandl2018}, which provides Fitzpatrick skin type annotations \cite{Fitzpatrick1988} enabling subgroup analysis. An EfficientNet-B0 model is pretrained on HAM10000 and fine-tuned on DDI. We then apply targeted synthetic data augmentation using color-space perturbations and geometric transforms, with augmentation factors swept from 1$\times$ to 10$\times$. Stages 2 and 3 share identical hyperparameters (lr\,$=$\,$5{\times}10^{-5}$, batch\,$=$\,32, pos\_weight\,$=$\,dynamic) for valid causal comparison. Synthetic images are generated exclusively from training-split dark-skin melanoma images ($n\,{=}\,29$, 10 variants each, 290 total). All experiments use seed 42 for reproducibility. Bootstrap 95\% confidence intervals (1{,}000 iterations) are reported throughout.

\textbf{Results:} The pretrained baseline exhibits limited discriminative ability on dark-skin melanoma detection (AUROC\,$=$\,0.58, sensitivity\,$=$\,0.0). Fine-tuning on DDI improves overall AUROC (0.64\,$\rightarrow$\,0.67) but widens the Light--Dark AUROC gap from $-0.05$ to $-0.33$ (Light\,$=$\,0.80\,[0.63--0.94] vs.\ Dark\,$=$\,0.47\,[0.26--0.70]). Adding 290 train-only synthetic dark melanomas does not significantly improve Dark AUROC (0.46\,[0.27--0.67], $p\,{=}\,0.524$ vs.\ finetuned). An ablation study over synthetic multiplier (0$\times$--10$\times$) reveals a non-monotonic relationship, with moderate augmentation showing marginal benefit and 10$\times$ degrading performance.

\textbf{Conclusions:} Simple photometric augmentation is insufficient to close the skin-tone performance gap on DDI with $n\,{=}\,29$ dark melanomas. Fine-tuning alone amplifies the Light--Dark AUROC gap ($-0.05$\,$\rightarrow$\,$-0.33$), demonstrating that naive transfer learning on small imbalanced data may worsen fairness. Adding synthetic augmentation from photometric transforms does not significantly improve Dark AUROC ($p\,{=}\,0.524$). This negative result highlights the need for more diverse generative models, larger underrepresented training sets, or fundamentally different mitigation strategies. Limitations include the small dataset size, single seed, and wide confidence intervals.
\end{abstract}
